\section{Discussion}
\label{sec:discussion}

\subsection{Why trace-grounding works}

The single sentence: \emph{every tool in a recorded trace was actually
invoked, so ``unnecessary tool'' is impossible by construction.} The
forward-distilled spec compiles a real trajectory into its effects,
then asks the candidate to \emph{produce those effects}, not
\emph{replay the trajectory}. RQ2 (\S\ref{sec:rq2}) quantifies the
structural property; RQ1 (\S\ref{sec:rq1}) quantifies what happens if
you forget it and revert to answer matching.

\subsection{What dynamism explains and what it doesn't}

RQ3's dominant \texttt{dynamism\_live} driver
(\S\ref{sec:rq3}) matches the intuition that ``the world moves between
distillation and evaluation.'' But dynamism is also collinear with
task content (\texttt{live\_read} tasks tend to be the
information-seeking ones); the path forward is the larger substrate
where content and dynamism can be partially decorrelated. We treat
the RQ3 result as the empirical defense of the \emph{stratification}
choice, not as a per-feature causal claim.

\subsection{Where the answer string is still useful}

DynamicMCPBench's refusal to grade the final answer is not a refusal to
\emph{report} it. The candidate's
\texttt{final\_assistant\_message} is stored, inspectable, and surfaced
in the report~--- it just doesn't move the score. This separation lets
the same artifacts feed downstream studies (e.g.\ fluency, justification
quality) without contaminating the trace score.

\subsection{Comparison shapes versus the headline}

The graph-sampling and direct-generation baselines in \S\ref{sec:rq2}
stay \emph{labeled} (\texttt{distiller\_version="baseline-..."}) at
every level~--- disk format, CLI, report aggregator, regression-tested
orthogonality guard. This was deliberate so the RQ2 comparison cannot
be confused for ``AGB without the graph''; the gap it measures is the
\emph{structural} one.
